\documentclass[10pt]{article}
\usepackage[T1]{fontenc}
\usepackage[utf8]{inputenc}
\usepackage[french]{babel}
\usepackage{amsmath,amssymb,mathrsfs,mathpazo}
\usepackage{geometry}
\geometry{a4paper,margin=2cm}
\usepackage{graphicx,array,tabularx,tabularray,colortbl,multirow,multicol,booktabs,enumitem,pifont,ulem,fancybox,diagbox,cancel}
\usepackage{pstricks,pstricks-add,pst-plot,pst-tree,pst-node,pst-text,pst-func,pst-math,pst-bspline,pst-3dplot,pst-eucl,pst-solides3d}
\let\pscalculate\undefined
\usepackage{pst-bezier}
\usepackage{tikz,pgfplots,tkz-tab}
\usetikzlibrary{arrows,arrows.meta,patterns,calc,positioning,trees,decorations.markings,shapes.geometric,shapes.misc,shapes.arrows,intersections,decorations.pathmorphing,decorations.pathreplacing,backgrounds,shadows,babel}
\pgfplotsset{compat=1.18}
\usepgfplotslibrary{fillbetween}
\usepackage[np]{numprint}
\usepackage{listings,makecell,pas-tableur,twoopt}
\let\barre\undefined
\usepackage[pstricks]{tablvar}
\definecolor{darkgreen}{RGB}{0,100,0}
\providecommand{\red}{\color{red}}
\providecommand{\blue}{\color{blue}}
\providecommand{\green}{\color{green}}
\providecommand{\black}{\color{black}}
\providecommand{\white}{\color{white}}
\providecommand{\gray}{\color{gray}}
\providecommand{\cyan}{\color{cyan}}
\providecommand{\magenta}{\color{magenta}}
\providecommand{\yellow}{\color{yellow}}
\providecommand{\violet}{\color{violet}}
\providecommand{\orange}{\color{orange}}
\providecommand{\euro}{\texteuro}
\providecommand{\EUR}{\texteuro}
\providecommand{\ieme}{\textsuperscript{e}}
\providecommand{\degres}{\ensuremath{^\circ}}
\providecommand{\degree}{\ensuremath{^\circ}}
\providecommand{\vv}[1]{\overrightarrow{#1}}
\providecommand{\Coord}[1]{#1}
\providecommand{\decosix}{}
\providecommand{\starredbullet}{\ensuremath{\bullet}}
\pagestyle{empty}
\setlength{\parindent}{0pt}
\usepackage{adjustbox}\DeclareUnicodeCharacter{2236}{:}\DeclareUnicodeCharacter{2190}{\ensuremath{\leftarrow}}
\providecommand{\R}{\mathbb{R}}
\providecommand{\N}{\mathbb{N}}
\providecommand{\Z}{\mathbb{Z}}
\providecommand{\Q}{\mathbb{Q}}
\providecommand{\D}{\mathbb{D}}
\providecommand{\C}{\mathbb{C}}
\providecommand{\e}{{\mathrm{e}}}
\providecommand{\vect}[1]{\overrightarrow{#1}}
\providecommand{\vectt}[1]{\overrightarrow{\mathrm{#1}}}
\providecommand{\vvt}[1]{\overrightarrow{\mathrm{#1}}}
\providecommand{\np}[2][]{{#2~#1}}
\providecommand{\vv}[1]{{\overrightarrow{#1}}}
\renewcommand{\d}{{\mathrm{d}}}
\providecommand{\degres}{{^\circ}}
\providecommand{\degre}{{^\circ}}
\providecommand{\blue}{\color{blue}}
\providecommand{\red}{\color{red}}
\providecommand{\enskip}{\quad}
\providecommand{\llbracket}{\lbrack\!\lbrack}
\providecommand{\rrbracket}{\rbrack\!\rbrack}
\providecommand{\euro}{{\text{€}}}
\providecommand{\no}{{\text{n°}}}
\providecommand{\interval}[2]{{[#1;#2]}}
\providecommand{\intervFO}[2]{{[#1;#2[}}
\providecommand{\dotfill}{\ldots}

\providecommand{\ds}{\displaystyle}
\providecommand{\textding}[1]{#1}
\providecommand{\Oij}{\(\mathrm{O};\vec{\imath},\vec{\jmath}\)}
\providecommand{\Oijk}{\(\mathrm{O};\vec{\imath},\vec{\jmath},\vec{k}\)}
\providecommand{\Ouv}{\(\mathrm{O};\vec{u},\vec{v}\)}
\providecommand{\intervFF}[2]{\ensuremath{[#1;#2]}}
\renewcommand{\euro}{\ifmmode\text{\texteuro}\else\texteuro\fi}
\DeclareUnicodeCharacter{00B0}{\ifmmode^{\circ}\else\textdegree\fi}
\DeclareUnicodeCharacter{20AC}{\ifmmode\text{\texteuro}\else\texteuro\fi}
% Énoncés : Éducation nationale. Transcription et corrigés : APMEP, Denis Vergès et contributeurs. Découpage et classement : CoopMaths, Sébastien Lozano, Éric Elter et contributeurs. Adaptation de la présentation : Maths Explorer.
% CC BY-SA : https://coopmaths.fr/www/about/
% https://coopmaths.fr/alea/static/bac/2025/tex/bac_2025_09_sujet1_asie_1_cor.tex
\begin{document}
\noindent\textit{APMEP / CoopMaths. CC BY-SA. Présentation adaptée par Maths Explorer.}\par\medskip
Soit $f$ la fonction définie sur $\R$ par $f(x)=x \e^{-2 x}$.

On admet que $f$ est deux fois dérivable sur $\R$ et on note $f'$ la dérivée de la fonction $f$.

On note $\mathcal{C}_{f}$ la courbe représentative de $f$ dans un repère orthonormé du plan.





\bigskip

\textbf{Affirmation 1.} Pour tout réel $x$, on a $f'(x)=(-2 x+1) \e^{-2 x}$.

$f(x)=x \e^{-2x}$ donc $f'(x)=1 \times \e^{-2x} + x \times (-2)\e^{-2x} = \left (-2x+1\right )\e^{-2x}$.

\hfill \textbf{Affirmation 1 vraie}

\bigskip

\textbf{Affirmation 2.} La fonction $f$ est une solution sur $\R$ de l'équation différentielle :
 $y'+2 y=\e^{-2 x}$.

$f'(x)+2f(x)= \left (-2x+1\right )\e^{-2x} +2x\e^{-2x} = -2x\e^{-2x} +\e^{-2x}+2x\e^{-2x}=\e^{-2x}$

\hfill \textbf{Affirmation 2 vraie}

\bigskip

\textbf{Affirmation 3.} La fonction $f$ est convexe sur $]-\infty ~;~ 1]$.

$f'(x)= \left (-2x+1\right )\e^{-2x}$ donc $f''(x)=(-2)\times \e^{-2x} + \left (-2x+1 \right )\times (-2)\e^{-2x} = \left (4x-4 \right ) \e^{-2x}$.

On étudie le signe de $f''(x)$ sur $\R$.

\[ \begin{tablvar}[intervalwidth=6em]{2}
 \hline
 x & -\infty & & 1 & & +\infty\\
 \hline
4x-4 & & - & \barre[0] & + & \\
 \hline
 \e^{-2x} & & + & \barre & + & \\
 \hline
f''(x) & & - & \barre[0] & + & \\
 \hline
f & & \text{concave} & \barre & \text{convexe} & \\
 \hline
 \end{tablvar}\]

\hfill \textbf{Affirmation 3 fausse}

	\bigskip

	\textbf{Affirmation 4.} L'équation $f(x)=-1$ admet une unique solution sur $\R$.

On détermine le signe de $f'(x) = \left (-2x+1\right )\e^{-2x}$ sur $\R$.

\[ \begin{tablvar}[intervalwidth=6em]{2}
 \hline
 x & -\infty & & \frac{1}{2} & & +\infty\\
 \hline
-2x+1 & & + & \barre[0] & - & \\
 \hline
 \e^{-2x} & & + & \barre & + & \\
 \hline
f'(x) & & + & \barre[0] & - & \\
 \hline
 \end{tablvar}\]

\begin{list}{\textbullet}{Valeurs remarquables}
\item $\ds\lim_{x\to -\infty} x =-\infty$ et $\ds\lim_{x\to -\infty} \e^{-2x}=+\infty$ donc $\ds\lim_{x\to -\infty} f(x)=-\infty$.
\item $f\left (\dfrac{1}{2}\right ) = \dfrac{1}{2} \e^{-2\times \frac{1}{2}} = \dfrac{1}{2}\e^{-1}>0$

\item $f(x)=x\e^{-2x}= \dfrac{1}{2}\times \dfrac{2x}{\e^{2x}}$; par croissance comparée, on a $\ds\lim_{x\to +\infty} \frac{2x}{\e^{2x}}=0$ donc $\ds\lim_{x\to +\infty} f(x)=0$.
\end{list}

On établit le tableau des variations de $f$ sur $\R$.

\[\begin{tablvar}[6em]{2}\hline
x & -\infty && \frac{1}{2} && +\infty \\\hline
f'(x) & & + & \barre[0] & - & \\\hline
\variations{\mil{f} & \bas{-\infty} && \haut{\frac{1}{2}\e^{-1}} &&  \bas{0}} \hline
\end{tablvar}\]

\begin{list}{\textbullet}{}
\item Sur l'intervalle $\left]-\infty\,;\,\frac{1}{2}\right]$, la fonction $f$ est continue et strictement croissante. De plus $\lim\limits_{x\to -\infty}f(x)=-\infty$ et $f\left(\frac{1}{2}\right)=\frac{1}{2}\e^{-1}>0$. Or $-1\in \left]-\infty\,;\,\frac{1}{2}\e^{-1}\right]$ donc, d'après le corollaire du théorème des valeurs intermédiaires, l'équation $f(x)=-1$ admet une unique solution sur $\left]-\infty\,;\,\frac{1}{2}\right]$.

\item Sur l'intervalle $\left[\frac{1}{2}\,;\,+\infty\right[$, la fonction $f$ est strictement positive donc l'équation $f(x)=-1$ n'admet aucune solution sur cet intervalle.
\end{list}

Donc l'équation $f(x)=-1$ admet une unique solution sur $\R$.

\hfill \textbf{Affirmation 4 vraie}

\bigskip

\textbf{Affirmation 5.} L'aire du domaine délimité par la courbe $\mathcal{C}_{f}$, l'axe des abscisses et les droites d'équation $x=0$ et $x=1$ est égale à $\dfrac{1}{4}-\dfrac{3 \e^{-2}}{4}$.

L'aire du domaine est égale à $\mathcal{A} = \ds\int_{0}^{1}f(x) \d x = \ds \int_{0}^{1} x \e^{-2x} \d x$.

On utilise une intégration par parties:
$\ds\int_{a}^{b} u(x)v'(x)\d x =\left [ u(x)v(x)\strut \right ]_{a}^{b} -\int_{a}^{b} u'(x) v(x) \d x$.

On pose
$\left \{
\begin{array}{l}
u(x)=x\\
v'(x)=\e^{-2x}
\end{array}
\right .$
donc
$\left \{
\begin{array}{l}
u'(x)=1\\
v(x)= -\dfrac{1}{2} \e^{-2x}
\end{array}
\right .$

$\begin{aligned}[t]
\mathcal{A}
& = \ds \int_{0}^{1} x \e^{-2x} \d x
 = \left [x\times \left ( -\dfrac{1}{2}\e^{-2x} \right ) \right ]_{0}^{1}- \ds\int_{0}^{1} 1\times \left ( -\dfrac{1}{2}\e^{-2x} \right ) \d x
 = \left [-\dfrac{x\e^{-2x} }{2} \right ]_{0}^{1}- \left [ \dfrac{1}{4}\e^{-2x} \right ]_{0}^{1}\\
& = \left [ - \dfrac{\e^{-2}}{2} - 0 \right  ] - \left [ \dfrac{\e^{-2}}{4}  - \dfrac{\e^{0}}{4}\right ]
= \dfrac{1}{4} - \dfrac{3\e^{-2}}{4}
\end{aligned}$

\hfill \textbf{Affirmation 5 vraie}

\bigskip

\end{document}
